\documentclass[11pt,a4paper]{article}
\usepackage[utf8]{inputenc}
\usepackage[T1]{fontenc}
\usepackage[margin=2.5cm]{geometry}
\usepackage{amsmath,amssymb,booktabs}
\setlength{\parindent}{0pt}
\setlength{\parskip}{6pt}

\title{Change request for MarkovSDDP4:\\ running rule for reserves}
\author{}
\date{26 November 2026}

\begin{document}
\maketitle

\section{Purpose}
Add one modelling feature to the existing weekly model in MarkovSDDP4
(\texttt{markov\_sddp.ipynb}, data in \texttt{duge.yaml}): a unit may only
supply reserve capacity when it is running, and it can only supply its full
reserve when it runs at or above its minimum power. The rule follows
Helseth, Fodstad and Mo (2016), \emph{IEEE Transactions on Sustainable
Energy} 7(3), equations (6) and (13), applied to each reserve direction and
to both the generating unit and the pumping unit.

Everything not listed in Section~3 stays exactly as it is.

\section{Existing notation used here}
\begin{tabular}{@{}ll@{}}
\toprule
Symbol & Meaning\\
\midrule
$t\in\mathcal T$ & hour within the week\\
$f\in\mathcal F^C$ & reserve product (currently only FCR-N)\\
$p^G_t$, $p^P_t$ & scheduled generation and pumping [MW]\\
$r^{G+}_{f,t}$ & reserve from increasing generation (up) [MW]\\
$r^{G-}_{f,t}$ & reserve from decreasing generation (down) [MW]\\
$r^{P+}_{f,t}$ & reserve from increasing pumping (down) [MW]\\
$r^{P-}_{f,t}$ & reserve from decreasing pumping (up) [MW]\\
$R^{G,\max}$, $R^{P,\max}$ & maximum reserve from each unit in each direction [MW]\\
\bottomrule
\end{tabular}

\section{What to change}

\subsection{New parameters}
The same parameters as in Helseth et al.\ (2016): a minimum power and a
factor $\gamma$ per unit.

\begin{tabular}{@{}lll@{}}
\toprule
Symbol & Meaning & Unit\\
\midrule
$P^{G,\min}$ & minimum output of the generating unit & MW\\
$P^{P,\min}$ & minimum pumping power of the pumping unit & MW\\
$\gamma^G$ & factor limiting the use of reserves from the generating unit & --\\
$\gamma^P$ & factor limiting the use of reserves from the pumping unit & --\\
\bottomrule
\end{tabular}

The factors are defined as in equation (13) of the paper:
\begin{equation}
\gamma^G=\max\left\{\frac{P^{G,\min}}{R^{G,\max}},\,1\right\},\qquad
\gamma^P=\max\left\{\frac{P^{P,\min}}{R^{P,\max}},\,1\right\}.
\tag{A0}
\end{equation}

Read $P^{G,\min}$ and $P^{P,\min}$ from \texttt{duge.yaml}, next to the
existing unit data and with the existing naming style. Do not hard-code
values and do not guess them: if a key is missing, stop with a clear error
message. The values are supplied by the user. Compute $\gamma^G$ and
$\gamma^P$ from (A0) when the data are loaded; they are not entries in the
data file.

No new variables, no new sets.

\subsection{Four constraints}
The constraints have the form of equation (6) of the paper,
$\gamma_h r_{kh}-\tfrac{1}{\tau_k}e_{kh}\le 0$, where
$\tfrac{1}{\tau_k}e_{kh}$ is the power of the unit, which is $p^G_t$ or
$p^P_t$ here. For every hour $t\in\mathcal T$, in every weekly subproblem
(every stage and every Markov state, forward and backward pass):
\begin{align}
\gamma^G\sum_{f\in\mathcal F^C} r^{G+}_{f,t}-p^G_t &\le 0
  \tag{A1}\\
\gamma^G\sum_{f\in\mathcal F^C} r^{G-}_{f,t}-p^G_t &\le 0
  \tag{A2}\\
\gamma^P\sum_{f\in\mathcal F^C} r^{P+}_{f,t}-p^P_t &\le 0
  \tag{A3}\\
\gamma^P\sum_{f\in\mathcal F^C} r^{P-}_{f,t}-p^P_t &\le 0
  \tag{A4}
\end{align}

\begin{itemize}
\item (A2) \textbf{replaces} the existing constraint
      $p^G_t-\sum_f r^{G-}_{f,t}\ge 0$.
\item (A4) \textbf{replaces} the existing constraint
      $p^P_t-\sum_f r^{P-}_{f,t}\ge 0$.
\item (A1) and (A3) are \textbf{new} and are added.
\end{itemize}

\subsection{Special case}
If $R^{G,\max}=0$ or $R^{P,\max}=0$, the reserves of that unit are already
zero and (A0) is undefined; set its factor to 1.

\section{What must not change}
\begin{itemize}
\item The objective function.
\item The upper unit limits $p^G_t+\sum_f r^{G+}_{f,t}\le P^{G,\max}$ and
      $p^P_t+\sum_f r^{P+}_{f,t}\le P^{P,\max}$.
\item The unit reserve limits ($\le R^{G,\max}$, $\le R^{P,\max}$), the
      commitment balances, the product limits and the symmetry requirement.
\item The water balance, the reservoir limits and the water margin for
      reserves.
\item The state variable (reservoir volume only), the cuts, the Markov
      chain, the end-of-horizon treatment, the training loop and the
      stopping rule.
\end{itemize}
The subproblem remains a linear programme. No binary variables.

\section{Checks after the change}
\begin{enumerate}
\item In every solved hour: if $p^G_t=0$ then
      $\sum_f r^{G+}_{f,t}=\sum_f r^{G-}_{f,t}=0$, and the same for the
      pumping unit.
\item In every solved hour: the generating unit supplies $R^{G,\max}$ in a
      direction only if $p^G_t\ge\gamma^G R^{G,\max}$, and the same for the
      pumping unit.
\item With the same data and seeds, the upper bound after training is not
      higher than before the change, because the feasible set has only
      become smaller. Report both values.
\item Report, before and after, the expected income from energy and from
      reserve capacity and the number of hours each unit runs, so the
      effect of the rule is visible.
\end{enumerate}

\end{document}
